\section[Nonstandard GCT: coefficient, spectrum and heat connection]{Nonstandard GCT: the coefficient, its arithmetic spectrum, and heat connection}
\label{sec:gct-arithmetic-spectrum}

\subsection{The original source and the direction of its programme}
Mulmuley's discussion of nonstandard Riemann hypotheses concerns a proposed
extension of the finite-field geometry behind canonical-basis positivity.
The actual maps in his account start with the plethysm representation
$H=GL_n(\C)\longrightarrow GL(V_\mu(H))$, its nonstandard quantum
deformation, and then the class-variety triple
$\bar H_{\hat h}\longrightarrow\bar H\longrightarrow GL(\bar W)$.
The intended destination is explicit representation-theoretic obstructions
to characteristic-zero complexity inclusions
\cite[section ``How to prove PH1 and why should it hold?'']{mulmuley2009gct}.
In that account, positivity supplies proposed polyhedral counting
descriptions; emptiness and eventual nonemptiness of the two different
polytopes supply the obstruction comparison. The source separately states
the positivity and obstruction hypotheses in its refined GCT6 theorem.
These are recorded as the author's research programme, not promoted here
to proved premises or reversed into an implication about classical RH.
The maps calculated below concern one explicit coefficient in the later
two-row construction, and the existing arithmetic spectral module.

Blasiak, Mulmuley and Sohoni's GCT IV prints a mixed-sign expression
\cite[\S17.2, example labelled ``ex not positive'']{blasiakMulmuleySohoni2013}.
Their parameter macro denotes $q$. Retain exactly their scalar expression
\begin{equation}\label{eq:gct-source-coefficient}
 \begin{split}
 C(q)&=\frac{[6]_q([7]_q-[3]_q)-[3]_q[8]_q}{[2]_q}\\
     &=q^{10}-q^4-q^{-4}+q^{-10}
       =(q-q^{-1})^2[7]_q[3]_q,\\
 [j]_q&=q^{j-1}+q^{j-3}+\cdots+q^{1-j}.
 \end{split}
\end{equation}
The first line is the source's printed computation with $n_3=3,n_2=2$.
Its literal target tableau has a weight mismatch, proved below; the
earlier identification of this expression as the coefficient for that
literal target is withdrawn. None of the polynomial, spectrum or
heat calculations below uses that identification.
Comparing its numerator with $[2]_q$ times the claimed Laurent polynomial
and expanding the finite sums proves the equality without division at
specializations. In particular the last two lines define the continued
coefficient even at $q=\pm1,\pm i$. At $q=\pm i$ the first line itself
has a zero denominator and is not an evaluation prescription.
The source's preceding theorem gives sign-definite Laurent coefficients
for the individual Chevalley generators; this divided-power example is
explicitly outside that stronger sign assertion. Its \S1.7 retains the
tableau rescaling $(-1/[2]_q)^{\deg T}$ and the integral form used at $q=1$.
Thus continuation of this single Laurent coefficient does not by itself
specialize the entire source basis at $[2]_q=0$.

\paragraph{The literal target and its exact weight projection.}
Write a column from top to bottom in parentheses, and put $p=[2]_q$.
Over $\Q(q)$ the source's displayed tableaux specialize to
\[
\begin{aligned}
 T&=(123)(123)(124)(12)(12)(1)(1),\\
 P&=-p^{-1}(123)(123)(124)(23)(12)(4)(1),\\
 T_1&=(123)(123)(124)(23)(12)(1)(1),\\
 T_2&=-p^{-1}(123)(123)(123)(12)(12)(4)(1).
\end{aligned}
\]
The source's alphabet is
$1\leftrightarrow(1,1)$, $2\leftrightarrow(1,2)$,
$3\leftrightarrow(2,1)$, $4\leftrightarrow(2,2)$.
Its nonstandard-column dictionary preserves the resulting counts:
for example $(23)=c^{VW}_{21,12}+p^{-1}c^{VW}_{21,21}$
has $V$-weight $(1,1)$ and $W$-weight $(1,1)$ in both summands,
and the two terms of $(124)$ both have weights $(2,1),(1,2)$.
Tensor products add weights; the source's based products and quotient
maps are $U_q(V)\otimes U_q(W)$-module maps and preserve them
\cite[nonstandard-column dictionary and Proposition labelled
``p NSC facts'']{blasiakMulmuleySohoni2013}. Counting gives
\[
\begin{array}{c|cc}
 &\operatorname{wt}_V&\operatorname{wt}_W\\\hline
 T&(12,3)&(9,6)\\
 T_1,T_2&(11,4)&(9,6)\\
 P&(10,5)&(8,7)\\
 Q=-p^{-1}(123)(123)(124)(13)(12)(4)(1)&(10,5)&(9,6).
\end{array}
\]
Here $Q$ records the explicit alternative endpoint, not a renaming
of $P$. On the weight decomposition $V_\nu=\bigoplus_{\alpha,\beta}
(V_\nu)_{\alpha,\beta}$ let $\pi_{\alpha,\beta}$ be the projection
which is the identity on the named summand and zero on the others.
Its kernel is the direct sum of all the other summands and its image
is exactly the named summand. Since $F_V$ has weight
$((-1,1),(0,0))$, the map
\[
 F_V^{(2)}:(V_\nu)_{(12,3),(9,6)}
       \longrightarrow(V_\nu)_{(10,5),(9,6)}
\]
is followed by zero under $\pi_{(10,5),(8,7)}$.
Thus the coefficient of the literal printed $P$ is zero.
This conclusion uses the generic source field, not a root-of-unity
specialization, and is unaffected by the nonzero rescaling $-p^{-1}$.
Replacing the indicated column $(23)$ by $(13)$ restores the target
weight and gives $Q$. Weight compatibility alone does not compute
its coefficient; its contextual straightening must be retained in
a two-step generator calculation. The subsequent scalar base change
applies to the explicit $C(q)$ in \eqref{eq:gct-source-coefficient}
independently of that calculation.

\subsection{A finite-free map retaining the original arithmetic coordinate}
Work first over a characteristic-zero field $k$. The original slice of
the polynomial inverse construction has $b=c=0$ and $r^2=-a$.
Adjoin $p$ by $p^2=4-a$. This gives an exact base change
\begin{equation}\label{eq:gct-conic-map}
 k[p]\otimes_{k[a]}k[a,r]/(r^2+a)
 =k[p,r]/(p^2-r^2-4)\simeq k[q,q^{-1}],
 \quad a\longmapsto4-p^2.
\end{equation}
The isomorphism and its inverse are
\[
 p=q+q^{-1},\quad r=q-q^{-1},\qquad
 q=(p+r)/2,\quad q^{-1}=(p-r)/2.
\]
Indeed the Laurent expressions have $p^2-r^2=4$, and in the conic
ring $(p+r)(p-r)=4$. The compositions fix all four generators, proving
both maps and their inverse identities. Retain $a=2s$, exactly as in
the preceding arithmetic construction. Then
\begin{equation}\label{eq:gct-s-map}
 s=s(q)=-\frac{(q-q^{-1})^2}{2}
       =1-\frac{q^2+q^{-2}}2.
\end{equation}
No original source coordinate is renamed: $p$ is the new cover coordinate,
whereas the original triple $(x,y,w)$ on the nonboundary chart is
\[
 x=-\frac1{2r},\quad y=3r,\quad w=26r^2,
 \qquad F(x,y,w)=(-r^2,0,0).
\]
These expressions are defined for $r\ne0$, equivalently $q\ne0,\pm1$.
The other original fibre point $(0,0,-r^2)$ remains present separately.

The Laurent ring is free of rank four over $k[s]$:
\begin{equation}\label{eq:gct-quartic}
 k[q,q^{-1}]\simeq k[s,q]/(q^4+(2s-2)q^2+1).
\end{equation}
To check this, the quartic yields
$q^{-1}=-q^3+(2-2s)q$; substitution of \eqref{eq:gct-s-map} gives the
inverse map. Division by the monic quartic gives a unique representative
of degree at most three. The second free basis is $(1,r,p,rp)$, because
$r^2=-2s$ and $p^2=4-2s$ are successive monic quadratic relations.
Explicit conversion is
\begin{gather*}
 r=q^3+(2s-1)q,\qquad p=-q^3+(3-2s)q,\qquad rp=2q^2+2s-2,\\
 q=(p+r)/2,\qquad q^2=(rp+2-2s)/2,\qquad
 q^3=r-(2s-1)(p+r)/2.
\end{gather*}
These identities also prove independence of the second basis directly.

For the actual topological module $(Q,S)$ of the preceding arithmetic
construction, take $k=\C$ and give
$Q^{(4)}=\C[q,q^{-1}]\otimes_{\C[s]}Q$ the fourfold direct-sum topology
in $(1,q,q^2,q^3)$. Here $s$ acts by the already defined continuous
operator $S$, not by a freely chosen replacement spectrum. Set
\begin{equation}\label{eq:gct-K}
 K=\begin{pmatrix}
 0&0&0&-I\\ I&0&0&0\\0&I&0&2I-2S\\0&0&I&0
 \end{pmatrix},\qquad \widehat S=\operatorname{diag}(S,S,S,S).
\end{equation}
The quartic proves
\begin{gather*}
 K^{-1}=-K^3+(2I-2\widehat S)K,\qquad
 \widehat S=-\tfrac12(K-K^{-1})^2,\\
 (K+K^{-1})^2=4I-2\widehat S.
\end{gather*}
All expressions are continuous. The constant-coefficient inclusion and
projection retain $S$ itself. This constructs the action of the Laurent
coefficient algebra; a representation of all the nonstandard quantum
group generators is not asserted by this tensor product.

\subsection{The complete resolvent and all spectral fibres}
Spectrum here means failure of a continuous two-sided inverse on the
specified topological space. No Banach-space spectral mapping theorem
is needed.

\begin{theorem}\label{thm:gct-resolvent}
For $\lambda\ne0$, put $s_\lambda=s(\lambda)$ and
$A_\lambda=S-s_\lambda I$. Define
\begin{align*}
 E_\lambda v&=(-\lambda^{-1}v,-\lambda^{-2}v,\lambda v,v),\\
 L_\lambda y&=\lambda^{-2}y_0+\lambda^{-1}y_1+y_2+\lambda y_3,\\
 H_\lambda y&=(-\lambda^{-1}y_0,
       -\lambda^{-2}y_0-\lambda^{-1}y_1,y_3,0),\\
 J_0v&=(0,0,v,0),\qquad P_3x=x_3.
\end{align*}
Then $K-\lambda I$ is continuously invertible exactly when $A_\lambda$
is continuously invertible. In that case
\begin{equation}\label{eq:gct-resolvent}
 (K-\lambda I)^{-1}
 =H_\lambda-\tfrac12E_\lambda A_\lambda^{-1}L_\lambda,
 \qquad
 A_\lambda^{-1}=P_3(K-\lambda I)^{-1}(-2J_0).
\end{equation}
Moreover
\begin{equation}\label{eq:gct-spectrum}
 \operatorname{Spec}K
 =s^{-1}(\operatorname{Spec}S)\subset\C^\times,
 \qquad
 \ker(K-\lambda I)=E_\lambda\ker(S-s_\lambda I),
\end{equation}
with inverse $P_3$ on the displayed kernels.
\end{theorem}
\begin{proof}
Writing $B=K-\lambda I$, its four equations $Bx=y$ give, successively,
\begin{align*}
 x_0&=-\lambda^{-1}(x_3+y_0),\\
 x_1&=-\lambda^{-2}x_3-\lambda^{-2}y_0-\lambda^{-1}y_1,\\
 x_2&=\lambda x_3+y_3,\\
 -2(S-s_\lambda I)x_3&=L_\lambda y.
\end{align*}
Conversely these equations imply all four rows of $Bx=y$ by substitution.
Thus the last equation determines $x_3$ uniquely when $A_\lambda$ is
invertible, and the other three determine $x$ by \eqref{eq:gct-resolvent}.
This inverse is continuous, since it is a finite sum of continuous maps.

In the reverse direction, direct multiplication gives
$BE_\lambda=-2J_0A_\lambda$ and $P_3E_\lambda=I$.
If $B$ is invertible, $A_\lambda v=0$ forces $v=0$.
For any $v$, solve $Bx=-2J_0v$; the last eliminated equation gives
$A_\lambda P_3x=v$. This proves surjectivity, the second inverse formula,
and its continuity. Taking $y=0$ in the same equations proves the kernel
isomorphism. Finally $K$ itself has the polynomial inverse already
displayed, so zero adds no point to its spectrum.
\end{proof}

\begin{proposition}\label{prop:gct-real-locus}
The complete inverse image of the real arithmetic axis is
\[
 \Gamma=\R^\times\ \cup\ \{q:|q|=1\}\ \cup\ i\R^\times.
\]
The respective images are $(-\infty,0]$, $[0,2]$, and $[2,\infty)$.
Consequently
\begin{equation}\label{eq:gct-real-spectrum}
 \operatorname{Spec}S\subset\R
 \quad\Longleftrightarrow\quad
 \operatorname{Spec}K\subset\Gamma.
\end{equation}
\end{proposition}
\begin{proof}
Write $q=\rho e^{i\theta}$ with $\rho>0$. Equation \eqref{eq:gct-s-map}
gives
$\operatorname{Im}s(q)=-(\rho^2-\rho^{-2})\sin(2\theta)/2$.
This is zero exactly when $\rho=1$ or $q$ lies on one of the two axes.
Substitution on the axes and circle gives the stated ranges.
For completeness, the full fibres, with independent signs
$\varepsilon,\delta\in\{1,-1\}$, are
\begin{align*}
 s<0:\quad&q=\tfrac12\bigl(\varepsilon\sqrt{4-2s}
                              +\delta\sqrt{-2s}\bigr),\\
 0<s<2:\quad&q=\tfrac12\bigl(\varepsilon\sqrt{4-2s}
                              +i\delta\sqrt{2s}\bigr),\\
 s>2:\quad&q=\tfrac i2\bigl(\varepsilon\sqrt{2s-4}
                              +\delta\sqrt{2s}\bigr).
\end{align*}
They follow by solving for $p,r$ in \eqref{eq:gct-conic-map}.
Their products $qq^{-1}=1$ ensure none is zero; away from $s=0,2$
they are distinct. At $s=0$ the quartic is $(q^2-1)^2$; at $s=2$
it is $(q^2+1)^2$. Every complex $s$ has a nonzero quartic root,
so \eqref{eq:gct-spectrum} proves both directions of
\eqref{eq:gct-real-spectrum}, including its converse.
\end{proof}

The two symmetries of each fibre are $q\mapsto-q$ and $q\mapsto q^{-1}$.
In the displayed $(p,r)$ signs, inversion changes $\delta$ and negation
changes both signs. The branch points are exactly $q=\pm1,\pm i$:
\[
 s'(q)=-q+q^{-3},\qquad s''(q)=-1-3q^{-4}=-4
 \quad(q^4=1).
\]
Thus their ramification index is two, not a removable multiplicity.
Indeed let $Sv=s_0v$ with $v\ne0$, $s_\lambda=s_0$ and
$\lambda^4=1$. The value $s_0=0$ is included. Then
\begin{equation}\label{eq:gct-Jordan}
 (K-\lambda I)E_\lambda v=0,\qquad
 (K-\lambda I)(\lambda^{-2}v,2\lambda^{-3}v,v,0)=E_\lambda v.
\end{equation}
To prove the second identity, differentiate
$(K-\lambda I)E_\lambda v=-2J_0(S-s_\lambda I)v$ in $\lambda$;
$s'_\lambda=0$ at the branch point. The two vectors are independent:
applying $K-\lambda I$ to a proposed dependence proves that its second
coefficient is zero. A semisimple source eigenvalue can therefore acquire
a genuine two-step Jordan chain on this ramified extension.

\subsection{Exact signs and three real structures}
The recurrence $U_0=1,U_1=p,U_j=pU_{j-1}-U_{j-2}$ gives
$U_j=[j+1]_q$ by cancellation of the interior Laurent terms. Hence
\begin{equation}\label{eq:gct-cs}
 C(q)=-a(7-14a+7a^2-a^3)(3-a)=c(s),\qquad
 c(s)=-2s(7-28s+28s^2-8s^3)(3-2s).
\end{equation}
In particular $C(K)=c(\widehat S)$. At $q=1$ the original vanishing
factor has $a=-(q-1)^2(q+1)^2/q^2$, so
$C(q)/(q-1)^2\to84$. In the original $a$ coordinate the zero is simple
with coefficient $-21$. Both orders are retained.

Let $\alpha_j=2\sin^2(j\pi/7)$, $j=1,2,3$. On the unit circle
$[7]_q=\sin(7\theta)/\sin\theta$, and therefore
\[
 c(s)=-32s(s-\alpha_1)(s-\alpha_2)(s-\tfrac32)(s-\alpha_3).
\]
To check the factorization, the cubic factor in \eqref{eq:gct-cs}
vanishes at the three distinct $\alpha_j$ and has leading coefficient
$-8$; multiplication by $-2s(3-2s)$ proves the formula. Since sine
is strictly increasing on $[0,\pi/2]$ and
$2\pi/7<\pi/3<3\pi/7<\pi/2$, the root order is
$0<\alpha_1<\alpha_2<3/2<\alpha_3<2$.
The signs of $c$ on the six successive complementary real intervals are
$+,-,+,-,+,-$. This proves $C(q)>0$ on real nonzero $q\ne\pm1$,
and $C(q)<0$ on every nonzero purely imaginary $q$; in particular
$C(\pm i)=-4$. The exact circle example
\[
 q_0=(\sqrt{14}+i\sqrt2)/4,\qquad s(q_0)=1/4,
 \qquad C(q_0)=-65/32
\]
follows by multiplication, with $q_0^{-1}=\overline{q_0}$.

Three conjugate-linear, multiplicative involutions on the Laurent
algebra extend the same involution $s^*=s$:
\[
 q^{*_{\rm r}}=q,\qquad q^{*_{\rm u}}=q^{-1},\qquad
 q^{*_{\rm i}}=-q.
\]
Each squares to the identity on $q$ and scalars, hence on every Laurent
polynomial. Their characters are respectively the real axis, unit circle
and imaginary axis, since the character condition is
$\overline{q_0}=q_0$, $q_0^{-1}$, or $-q_0$.
They fix $s$ by \eqref{eq:gct-s-map}. This is an explicit relation
between the three real forms, not an identification of their positivity
cones. For example, in a positive Hilbert-space representation with
bounded invertible $K$, the first involution gives
$\widehat S=-(K-K^{-1})^2/2\le0$. The second makes $K$ unitary and gives
$0\le\widehat S=I-\operatorname{Re}(K^2)\le2I$.
For the third write $K=iA$ with bounded self-adjoint invertible $A$.
Then $\widehat S=(A+A^{-1})^2/2\ge2I$, since
$(A+A^{-1})^2-4I=(A-A^{-1})^2\ge0$.
These are statements about those specified representations; they do
not assume a positive Hilbert realization of the arithmetic quotient.

The algebraic fibre trace of multiplication by $C$ is $4c(s)$ even at
ramification, because the module is free of rank four and $C=c(s)$.
It is not the quadratic Weil value of a scalar test. The precise test
map explains the distinction. If $h$ is a scalar transform in the
retained test space, multiplication sends it to $c(s)h(s)$; the reflected
sesquilinear spectral summand is consequently multiplied by
$c(s)\overline{c(\bar s)}$. At real $s$ this is $|c(s)|^2$, not $c(s)$.
Equivalently, for an $S$-symmetric sesquilinear form $B$ on an invariant
polynomial domain,
\[
 B(c(S)u,c(S)u)=B(u,c(S)^2u),
\]
because $c$ has real coefficients and the adjoint can be moved one
factor at a time. An inserted value $B(u,c(S)u)$ and the actual value
$B(c(S)u,c(S)u)$ are thus related by different, explicit operators.
The negative coefficient above supplies no negative scalar Weil value
by this multiplication map. A spectral point of $K$ outside $\Gamma$,
on the other hand, would transfer exactly by
Theorem~\ref{thm:gct-resolvent} to an off-real arithmetic spectral point.

\subsection{All four material and arithmetic heat sectors}
Put $B_0=\C[s,s^{-1},(s-2)^{-1}]$ and
\[
 \mathcal A=B_0[r,p]/(r^2+2s,\ p^2+2s-4).
\]
This is the same cover localized away from its two branch values.
There is a unique derivation $D$ extending $\partial_s$ on $B_0$:
\begin{equation}\label{eq:gct-derivation}
 Dr=-1/r=r/(2s),\qquad Dp=-1/p=p/(2(s-2)),
 \qquad Dq=-q/(rp).
\end{equation}
Differentiating the two defining relations forces the first two
identities, because $r,p$ are invertible. Conversely these assignments
annihilate both relations, so they define a derivation. The conic map
then gives the last identity. Thus $D$ is the original material
coefficient derivative with $a'=2,b=c=0$, and $Ds=1$.

\begin{proposition}\label{prop:gct-four-heat}
For a section $f=f_0+r f_1+p f_2+rp f_3$ of this module, the coefficient
vector of $Df$ is $(\partial_s+A)\mathbf f$, where
\begin{equation}\label{eq:gct-four-connection}
 A=\operatorname{diag}\left(0,\frac1{2s},\frac1{2(s-2)},
                          \frac1{2s}+\frac1{2(s-2)}\right).
\end{equation}
Its full square, including its zeroth-order terms, is
\begin{equation}\label{eq:gct-four-heat}
 D^2=\partial_s^2+2A\partial_s+
 \operatorname{diag}\left(0,-\frac1{4s^2},-
    \frac1{4(s-2)^2},-\frac1{s^2(s-2)^2}\right).
\end{equation}
The lifted arithmetic heat generator is $-D^2/4$, retaining the
original sign and factor. The field trace is $\operatorname{Tr}f=4f_0$,
with kernel $rB_0\oplus pB_0\oplus rpB_0$, and
$\operatorname{Tr}Df=\partial_s\operatorname{Tr}f$.
\end{proposition}
\begin{proof}
Apply the product rule and \eqref{eq:gct-derivation} to each term of $f$.
Squaring $\partial_s+A$ gives
$\partial_s^2+2A\partial_s+A'+A^2$. The middle entries of $A'+A^2$
are $-1/(4s^2)$ and $-1/(4(s-2)^2)$. The last one is
\[
 -\frac1{4s^2}-\frac1{4(s-2)^2}+\frac1{2s(s-2)}
 =-\frac1{s^2(s-2)^2},
\]
proving every lower-order term. The four sign automorphisms
$(r,p)\mapsto(\delta r,\varepsilon p)$ have sum $4f_0$; alternatively
the multiplication matrices of $r,p,rp$ have zero diagonal, giving
the trace even before localization. Linear independence of the free
basis identifies the kernel. The first row of $A$ is zero, proving
the trace-derivative identity without discarding the other three sectors.
\end{proof}

On a simply connected analytic neighbourhood avoiding $0,2$, choose
nonvanishing branches $r(s),p(s)$. The four evaluations are
\[
 h_{\delta,\varepsilon}(s)=f_0(s)+\delta r(s)f_1(s)
       +\varepsilon p(s)f_2(s)+\delta\varepsilon r(s)p(s)f_3(s).
\]
Their inverse recovers $f_0$ by averaging, $f_1$ by averaging
$\delta h/r$, $f_2$ by averaging $\varepsilon h/p$, and $f_3$ by
averaging $\delta\varepsilon h/(rp)$, each with factor $1/4$.
These formulas prove an invertible local map and show directly that
evaluation intertwines $D$ with $\partial_s$ and $D^2$ with
$\partial_s^2$. Equivalently, with
$E=\operatorname{diag}(1,r,p,rp)$,
$\partial_s+A=E^{-1}\partial_s E$ on coefficient vectors.
This does not erase the factors: the inverse ceases to exist at a branch
value. Around $s=0$, $r$ and $rp$ change sign; around $s=2$, $p$ and
$rp$ change sign. The residues of $A$ are respectively
$\operatorname{diag}(0,1/2,0,1/2)$ and
$\operatorname{diag}(0,0,1/2,1/2)$, retaining exactly this monodromy.

For a finite material increment $\tau$, local continuation gives
\[
 s_\tau=s+\tau,\quad r_\tau^2=r^2-2\tau,\quad
 p_\tau^2=p^2-2\tau,\quad q_\tau=(p_\tau+r_\tau)/2.
\]
The initial branches fix $r_0=r,p_0=p$ and their continuation until a
branch value is met. Since $p_\tau^2-r_\tau^2=4$, this preserves the
conic and $q_\tau^{-1}=(p_\tau-r_\tau)/2$.
Differentiation recovers \eqref{eq:gct-derivation}. The original source
triple and its velocity are therefore
\[
 (-1/(2r_\tau),3r_\tau,26r_\tau^2),\qquad
 (-1/(2r_\tau^3),-3/r_\tau,-52).
\]
At $r_\tau=0$ its first coordinate escapes; at $p_\tau=0$ that source
triple remains finite and only the additional conic coordinate is
ramified. The finite boundary section $(0,0,2(s+\tau))$ has velocity
$(0,0,2)$.

These derivations are on the displayed localized coefficient module and
its analytic sections. On the actual arithmetic quotient, multiplication
has already been constructed by \eqref{eq:gct-K}, while differentiation
uses the transported image and topology established in the preceding
heat section. For scalar initial data the present formulas give exactly
the pullback of $\Xi(s)=\xi(1/2+is)=8H_0(2s)$ and of its established
heat family, through \eqref{eq:gct-s-map}; no new initial function,
discarded coefficient factor, or substituted spectral coordinate occurs.

All ring, resolvent, ramification, coefficient and connection identities
have exact reproducible symbolic checks accompanying these full proofs.
The spectral criterion is an exact reduction, not a constructed point
outside $\Gamma$ or a complexity separation theorem.
